
\setcounter{footnote}{120}

\medskip
\noindent\SGAInlineRefTarget{sga3:iii:paragraph:4-23}\textbf{4.23. ---}
Still under the hypotheses of \SGARefLink{sga3:iii:theorem:4-21}{4.21}, we shall now study how the set of
subgroups \(Y\) lifting \(Y_{\mathscr J}\) behaves under conjugation by
sections of \(X\). If \(x\) is a section of \(X\) over \(S\) inducing the
unit section of \(X_{\mathscr J}\), the inner automorphism
\(\operatorname{Int}(x)\) defined by \(x\) sends flat subgroups of \(X\)
lifting \(Y_{\mathscr J}\) to flat subgroups of \(X\) lifting
\(Y_{\mathscr J}\). Under the conditions of \SGARefLink{sga3:iii:theorem:4-21:item:ii}{4.21(ii)}, the set of these
subgroups is a principal homogeneous space under \(Z^1(Y_0,N_0)\). We
shall see that there is then a subgroup
\(\Gamma\) of \(B^1(Y_0,N_0)\)\footnote{Editor's note: \(Z^1(Y_0,N_0)\)
has been replaced by \(B^1(Y_0,N_0)\), since the proof shows that
\(\Gamma\) is a subgroup of \(B^1(Y_0,N_0)\); cf. \SGARefLink{sga3:iii:lemma:4-27}{4.27}--\SGARefLink{sga3:iii:corollary:4-29}{4.29}.} such that
two subgroups of \(X\), flat over \(S\) and lifting \(Y_{\mathscr J}\),
are conjugate by an \(x\in X(S)\) inducing the unit of
\(X(S_{\mathscr J})\) if and only if their \emph{difference} in
\(Z^1(Y_0,N_0)\) belongs to \(\Gamma\). In the best cases we shall show
that \(\Gamma=B^1(Y_0,N_0)\). Thus the set of flat subgroups of \(X\)
lifting \(Y_{\mathscr J}\), modulo conjugation by the elements
\(x\in X(S)\) inducing the unit of \(X(S_{\mathscr J})\), is either empty
or a principal homogeneous space under \(H^1(Y_0,N_0)\) (cf. \SGARefLink{sga3:iii:corollary:4-29}{4.29} and
\SGARefLink{sga3:iii:proposition:4-36}{4.36}).

\medskip
\noindent\SGAInlineRefTarget{sga3:iii:paragraph:4-24}\textbf{4.24. ---}
We retain the notation of \SGARefLink{sga3:iii:theorem:4-21}{4.21}. Let \(Y\) be a flat subgroup of \(X\)
whose reduction is \(Y_{\mathscr J}\). Recall that in \SGARefLink{sga3:iii:remark:0-5}{0.5} we introduced
the functor \(L_0^X\) (respectively \(L_0^Y\)) defined, for a variable
\(S_0\)-scheme \(T\), by the identity
\[
 \operatorname{Hom}_{S_0}(T,L_0^X)
 =\operatorname{Hom}_{\mathscr O_T}
 \left(
  \omega^1_{X_0/S_0}\otimes_{\mathscr O_{S_0}}\mathscr O_T,
  \mathscr J\otimes_{\mathscr O_{S_0}}\mathscr O_T
 \right),
\]
and similarly with \(Y\) in place of \(X\), as well as the functor
\[
 L'_X=\prod_{S_0/S}L_0^X.
\]
We then have:

\medskip
\noindent\SGAInlineRefTarget{sga3:iii:lemma:4-25}\textbf{Lemma 4.25.}
\textit{There is a canonical exact sequence of
\(Y_0\)-\(\mathscr O_{S_0}\)-modules}
\[
 \tag{+}\label{sga3:iii:lemma:4-25:equation:plus}
 \mathbf n_{Y_0/X_0}
 \xrightarrow{d}
 \omega^1_{X_0/S_0}
 \longrightarrow
 \omega^1_{Y_0/S_0}
 \longrightarrow 0
\]
\textit{having the following properties:}
\begin{enumerate}
\item[\SGAInlineRefTarget{sga3:iii:lemma:4-25:item:i}(i)] \textit{After pullback to \(Y_0\), \(d\) gives the morphism
  \(\overline D_0\) of \SGARefLink{sga3:iii:proposition:4-8:item:iii}{4.8(iii)}.}
\item[\SGAInlineRefTarget{sga3:iii:lemma:4-25:item:ii}(ii)] \textit{If \(X_0\) and \(Y_0\) are smooth over \(S_0\), then
  \(d\) is injective. Since the two modules of differentials are then
  locally free of finite type, the same is true of
  \(\mathbf n_{Y_0/X_0}\), and the sequence is locally split.}
\end{enumerate}

\textit{Proof.}\footnote{Editor's note: The proof has been expanded to
take account of the additions made in Expos\'e I, Section 6.8.}
Let \(\pi_0:Y_0\to S_0\) denote the structure morphism. By SGA 1 II,
formula (4.3) (see also EGA IV\textsubscript{4}, 16.4.21), there is a
canonical exact sequence of \(\mathscr O_{Y_0}\)-modules
\[
 \tag{\(\dagger\)}\label{sga3:iii:lemma:4-25:equation:dagger}
 \mathscr N_{Y_0/X_0}
 \xrightarrow{\overline D_0}
 \Omega^1_{X_0/S_0}\otimes_{\mathscr O_{X_0}}\mathscr O_{Y_0}
 \longrightarrow
 \Omega^1_{Y_0/S_0}
 \longrightarrow 0.
\]
Since this sequence consists of
\((Y_0\mathbin{\times_S}Y_0)\)-equivariant modules and morphisms, its
pullback \((+)\) by \(\varepsilon_0^*\) is an exact sequence of
\(Y_0\)-\(\mathscr O_{S_0}\)-modules, and \((\dagger)\) is the pullback
of \((+)\) by \(\pi_0^*\) (cf. \SGARefLink{sga3:i:subsection:6-8}{Expos\'e I, Section 6.8}). This proves \SGARefLink{sga3:iii:lemma:4-25:item:i}{(i)}.

Suppose in addition that \(X_0\) and \(Y_0\) are smooth over \(S_0\).
Then, by SGA 1 II, 4.10 (see also EGA IV\textsubscript{4}, 17.2.3(i)
and 17.2.5), \(\overline D_0\) is injective and the sequence
\((\dagger)\) consists of \(\mathscr O_{Y_0}\)-modules locally free of
finite type (and is therefore locally split). By the equivalence of
categories I, \SGARefLink{sga3:i:proposition:6-8-1}{6.8.1}, \(d\) is also injective, and consequently the
sequence \((+)\) has the asserted properties.

\medskip
\noindent\SGAInlineRefTarget{sga3:iii:paragraph:4-26}\textbf{4.26. ---}\footnote{Editor's note: The original has
been expanded to make clear that this is an exact sequence of
\(Y_0\)-\(\mathscr O_{S_0}\)-modules.}
For every \(S_0\)-scheme \(f:T\to S_0\), the sequence \((+)\) gives an
exact sequence of \(Y_0(T)\)-\(\mathscr O(T)\)-modules
\[
\begin{aligned}
0&\longrightarrow
 \operatorname{Hom}_{\mathscr O_T}
 \left(f^*\omega^1_{Y_0/S_0},f^*\mathscr J\right)
 \longrightarrow
 \operatorname{Hom}_{\mathscr O_T}
 \left(f^*\omega^1_{X_0/S_0},f^*\mathscr J\right)\\
&\hspace{10em}\longrightarrow
 \operatorname{Hom}_{\mathscr O_T}
 \left(f^*\mathbf n_{Y_0/X_0},f^*\mathscr J\right).
\end{aligned}
\]
Thus there is an exact sequence of
\(Y_0\)-\(\mathscr O_{S_0}\)-modules
\[
 \tag{4.26.1}\label{sga3:iii:paragraph:4-26:equation:4-26-1}
 0\longrightarrow L_0^Y\longrightarrow L_0^X
 \xrightarrow{d}N_0.
\]
It gives an exact sequence of complexes of abelian groups
\[
 0\longrightarrow C^*(Y_0,L_0^Y)
 \longrightarrow C^*(Y_0,L_0^X)
 \xrightarrow{d^*}C^*(Y_0,N_0),
\]
and, in particular, the following commutative diagram with exact rows:

\begin{center}
 \SGAThreeNativeInput[0.88\textwidth]
   {native_diagrams/exp3/diagram_060_exp3_p069_exact_complex_rows.tex}
\end{center}

Observe that \(C^0(Y_0,L_0^Y)\) (respectively \(C^0(Y_0,L_0^X)\)) is
just
\[
 \operatorname{Hom}_{S_0}(S_0,L_0^Y)
 =\operatorname{Hom}_{S}(S,L'_Y)
\]
(respectively the analogous group with \(X\) in place of \(Y\)); by
\SGARefLink{sga3:iii:corollary:0-9}{0.9}, it is the group of sections of \(Y\) (respectively \(X\)) over
\(S\) inducing the unit section of \(X_{\mathscr J}\). Also,
\(d^1\) is precisely the morphism \(v_{i_{Y_0}}\) of \SGARefLink{sga3:iii:proposition:4-8:item:iii}{4.8(iii)}, where
\(i_{Y_0}:Y_0\to X_0\) is the canonical immersion.\footnote{Editor's
note: This follows from the definition of
\(d:\mathbf n_{Y_0/X_0}\to\Omega^1_{X_0/S_0}\) (cf. \SGARefLink{sga3:iii:lemma:4-25}{4.25}) and from the
definition of \(v_{i_{Y_0}}\) (cf. \SGARefLink{sga3:iii:proposition:4-8}{4.8}).}

\medskip
\noindent\SGAInlineRefTarget{sga3:iii:lemma:4-27}\textbf{Lemma 4.27.}
\textit{Under the conditions of \SGARefLink{sga3:iii:theorem:4-21}{4.21} for \(S\), \(\mathscr I\),
\(\mathscr J\), and \(X\), let \(Y\) be a subgroup of \(X\), flat over
\(S\) and lifting \(Y_{\mathscr J}\). Let \(i:Y\hookrightarrow X\)
denote the canonical immersion.}\footnote{Editor's note: Part \SGARefLink{sga3:iii:lemma:4-27:item:i}{(i)} below
has been added; it will be useful in \SGARefLink{sga3:iii:proposition:4-35-1}{4.35.1} and then in \SGARefLink{sga3:iii:paragraph:4-38:item:4}{4.38(4)} and \SGARefLink{sga3:iii:paragraph:4-20:equation:5}{(5)}.}
\begin{enumerate}
\item[\SGAInlineRefTarget{sga3:iii:lemma:4-27:item:i}(i)] \textit{Let \(i':Y\to X\) be a morphism of \(S\)-schemes
  lifting \(i_0\) (so that \(i'\) is also an immersion), let
  \(Y'=i'(Y)\), and let \(d(i,i')\) be the element of
  \(C^1(Y_0,L_0^X)\) such that \(i'=d(i,i')\cdot i\) (cf. \SGARefLink{sga3:iii:proposition:1-2-4}{1.2.4}).
  Then the deviation \(d(Y,Y')\in C^1(Y_0,N_0)\) (cf. \SGARefLink{sga3:iii:remark:4-5-1}{4.5.1}) is given
  by}
  \[
   d(Y,Y')=d^1\bigl(d(i,i')\bigr).
  \]
\item[\SGAInlineRefTarget{sga3:iii:lemma:4-27:item:ii}(ii)] \textit{Let \(x\in C^0(Y_0,L_0^X)\) be a section of \(X\)
  over \(S\) inducing the unit section of \(X_{\mathscr J}\) over
  \(S_{\mathscr J}\). Then the deviation
  \(d(Y,\operatorname{Int}(x)Y)\in C^1(Y_0,N_0)\) (cf. \SGARefLink{sga3:iii:remark:4-5-1}{4.5.1}) is given
  by}
  \[
   -d\bigl(Y,\operatorname{Int}(x)Y\bigr)
   =d^1\partial x=\partial d^0x.
  \]
\end{enumerate}

Indeed, \(Y'\) is the image of \(Y\) under the composite morphism
\footnote{Editor's note: Recall that
\(L'_X=\prod_{S_0/S}L_0^X\).}
\[
 Y\xrightarrow{(d(i,i'),i)}L'_X\mathbin{\times_S}X
 \longrightarrow X,
\]
which is denoted \(d(i,i')\cdot i\) in \SGARefLink{sga3:iii:proposition:4-8:item:iii}{4.8(iii)}. By loc. cit. and the
equality \(v_{i_0}=d^1\), we therefore have
\[
 \begin{aligned}
 &c\bigl(X,Y',d(i,i')\cdot i\bigr)-c(X,Y',i)\\
 &\qquad=v_{i_0}\bigl(d(i,i')\bigr)
 =d^1\bigl(d(i,i')\bigr).
 \end{aligned}
\]
But \(d(i,i')\cdot i=i'\) factors through \(Y'\) by definition, so the
first term is zero. Moreover, by \SGARefLink{sga3:iii:proposition:4-8:item:ii}{4.8(ii)},
\[
 c(X,Y',i)=d(Y',Y)=-d(Y,Y').
\]
Hence \(d(Y,Y')=d^1(d(i,i'))\), which proves \SGARefLink{sga3:iii:lemma:4-27:item:i}{(i)}.

Now let \(x\) be as in \SGARefLink{sga3:iii:lemma:4-27:item:ii}{(ii)}. From the formula
\[
 xyx^{-1}=xyx^{-1}y^{-1}y
 =(x-\operatorname{Ad}(y)x)y=(-\partial x)(y)\cdot y,
\]
we see that \(Y'\) is the image of \(Y\) under the immersion
\(i'=(-\partial x)\cdot i_Y\). Part \SGARefLink{sga3:iii:lemma:4-27:item:i}{(i)} therefore gives
\[
 -d\bigl(Y,\operatorname{Int}(x)Y\bigr)
 =d^1\partial x=\partial d^0x.
\]

\medskip
\noindent\SGAInlineRefTarget{sga3:iii:corollary:4-28}\textbf{Corollary 4.28.}
\textit{Two subgroups \(Y\) and \(Y'\) of \(X\), flat over \(S\) and
lifting \(Y_{\mathscr J}\), are conjugate by a section of \(X\) over
\(S\) inducing the unit section of \(X_{\mathscr J}\) if and only if}
\[
 d(Y,Y')\in
 \partial d^0 C^0(Y_0,L_0^X)
 \subset \partial C^0(Y_0,N_0)
 =B^1(Y_0,N_0).
\]
\noindent\emph{Source note.} The French re-edition prints \(G\) here;
the ambient group throughout this section is \(X\).

\medskip
\noindent\SGAInlineRefTarget{sga3:iii:corollary:4-29}\textbf{Corollary 4.29.}
\textit{If \(d^0\) is surjective, then \(Y\) and \(Y'\) as above are
conjugate by a section of \(X\) over \(S\) inducing the unit section of
\(X_{\mathscr J}\) if and only if \(d(Y,Y')\in B^1(Y_0,N_0)\).}

\medskip
\noindent\SGAInlineRefTarget{sga3:iii:corollary:4-30}\textbf{Corollary 4.30.}
\textit{Let \(Y\) be as in \SGARefLink{sga3:iii:lemma:4-27}{4.27}. The set of conjugates of \(Y\) by
sections of \(X\) over \(S\) inducing the unit section of
\(X_{\mathscr J}\) is isomorphic to}
\[
 d^1\partial C^0(Y_0,L_0^X)
 =C^0(Y_0,L_0^X)\big/\ker(d^1\partial).
\]

Observe now that
\(C^0(Y_0,L_0^X)/\ker(d^1\partial)\) is computed using only the left-hand
square of the commutative diagram in \SGARefLink{sga3:iii:paragraph:4-26}{4.26}. In particular, it can also
be computed in any diagram of the same type having the same left-hand
square. Consider the functor \(L_0^X/L_0^Y\) over \(S_0\), defined by
\[
 \operatorname{Hom}_{S_0}(T,L_0^X/L_0^Y)
 =\operatorname{Hom}_{S_0}(T,L_0^X)
  \big/\operatorname{Hom}_{S_0}(T,L_0^Y).
\]
There is a commutative diagram

\begin{center}
 \SGAThreeNativeInput[0.91\textwidth]
   {native_diagrams/exp3/diagram_061_exp3_p070_quotient_exact_complex_rows.tex}
\end{center}

whence, by the preceding observation:

